
On
poursuit la démonstration par récurrence sur le rang $r$ de $P$. Supposons le résultat prouvé pour les torseurs sous des tores déployés de rang $<r$. On écrit alors $M=M_1\oplus M_2$, avec $M_1\simeq\mathbb Z^{r-1}$ et $M_2\simeq\mathbb Z$; on pose $\underline{A}'=\bigoplus_{m\in M_1}\underline{A}_m$ et $\underline{A}''=\bigoplus_{m\in M_2}\underline{A}_m$, de sorte que $\underline{A}=\underline{A}'\otimes_{\underline{O}_S}\underline{A}''$. Soit $P_1=\Spec(\underline{A}')$; alors $P_1$ est un torseur sous $D_S(M_1)$ et $P$ est un torseur sur $P_1$, de groupe $\underline{G}_{m\,P_1}$ et d'algèbre l'image inverse de $\underline{A}''$ sur $P_1$. D'après ce qui précède, $K_\bullet(P)$ est le quotient de $K_\bullet(P_1)$ par le sous-groupe engendré par les $(\eta(m_2)-1)K_\bullet(P_1)$ pour $m_2\in M_2$. Par hypothèse de récurrence, $K_\bullet(P_1)$ est le quotient de $K_\bullet(S)$ par le sous-groupe engendré par les $(\eta(m_1)-1)K_\bullet(S)$ pour $m_1\in M_1$. Compte tenu de la relation
\[
\eta(m_1+m_2)-1=\eta(m_1)\eta(m_2)-1
=\eta(m_1)(\eta(m_2)-1)+\eta(m_1)-1,
\]
$K_\bullet(P)$ est donc le quotient de $K_\bullet(S)$ par le sous-groupe engendré par les $(\eta(m)-1)K_\bullet(S)$ pour $m\in M$, ou encore pour $m$ parcourant un système de générateurs de $M$.

\textbf{Corollaire 2.3.} \uline{Avec les notations et les hypothèses de} 2.1, \uline{supposons que} $S$ \uline{soit régulier}; \uline{alors on a un isomorphisme canonique} :
\[
K^{\bullet}(P)\simeq K^{\bullet}(S)/\bigl((\eta(m)-1)_{m\in M}\bigr)
\]
\uline{où} $((\eta(m)-1)_{m\in M})$ \uline{est l'idéal de} $K^{\bullet}(S)$ \uline{engendré par les} $\eta(m)-1$ \uline{pour} $m$ \uline{parcourant un système de générateurs de} $M$.

Si $S$ est régulier, l'homomorphisme canonique $K^{\bullet}(S)\to K_\bullet(S)$ est un isomorphisme (IV 2.5). De même, $P$ étant alors régulier, $K^{\bullet}(P)\to K_\bullet(P)$ est un isomorphisme. Le corollaire résulte donc immédiatement de 2.2.

